\section{Research context}
\label{sec:context}

\subsection{Host institution and organisational structure}

Forschungszentrum J\"ulich (FZJ), a member of the Helmholtz Association, is one of
Europe's largest interdisciplinary research centres; its work on energy, information and
bioeconomy rests on physics, materials science, neuroscience and, through the J\"ulich
Supercomputing Centre (JSC), large-scale simulation and high-performance computing. The
internship took place at IAS-7 (Civil Safety Research) of the Institute for Advanced
Simulation, directed by Prof.\ Dr.\ Armin Seyfried. IAS-7 pursues basic research on
pedestrian and fire dynamics through experiments, modelling and computer-based
simulation, organised in four divisions: Fire Dynamics, Pedestrian Dynamics --
Empiricism, Pedestrian Dynamics -- Modelling, and Pedestrian Dynamics -- Social
Psychology. Postdoctoral and doctoral researchers, research engineers, technical staff
who run the experimental series, and students work across these divisions, and the
institute develops and maintains open tools for the community: PeTrack~\cite{petrack},
which extracts pedestrian trajectories from video, and the simulator
JuPedSim~\cite{jupedsim}, alongside a public archive of experimental data. Current
projects include work on 3D motion in crowds and on pushing in dense crowds.

\subsection{Unified 3D body models in pedestrian dynamics}

A trajectory is a single point per person per frame, yet the phenomena studied in dense
crowds are properties of bodies. Recent empirical work at the institute used pressure
sensors and overhead cameras with 3D motion capture to follow how controlled forward
pushes propagate through standing crowds, and found that motion is detected more
sensitively from the velocity of three body landmarks than from the displacement of
the centre of mass~\cite{feldmann}; recent modelling couples the dynamics of the upper
body and the legs to explain emergent waves and rotation in ultra-dense crowds, where
physical contact is unavoidable~\cite{chatagnon}. A unified 3D body model across the
divisions would serve two uses. For \emph{inference}, it would extract movement
patterns by fusing computer vision with inertial suit motion capture where available.
For \emph{analysis}, it would support checks of contacts and collisions, to
investigate force propagation, and of visibility under different body attributes. The
aim of this work for that setting is transferable knowledge: understanding how to
build new parametric models from the ground up, with a toolset that is body-plan
agnostic and therefore not locked to one body model.

\subsection{From SMPL to a body-plan-agnostic pipeline}

The parametric-model lineage SMPL~\cite{smpl}, SMAL~\cite{smal} and
SMILify~\cite{smilify} is the starting point. SMPL represents the body as a rest-pose
template deformed by identity-dependent blend shapes and pose-dependent corrective blend
shapes, skinned linearly to a skeleton whose joints are regressed from the vertices, and
was learned from thousands of aligned 3D scans~\cite{smplweb}; SMAL relaxed the
fixed-topology and cooperative-subject assumptions for quadrupeds, and SMILify accepts an
arbitrary rigged template. Three properties of this setting make a from-scratch pipeline
valuable. First, widely used body models are distributed under registration and licence
agreements with separate commercial terms~\cite{smplweb}, whereas a pipeline that builds
shape and pose spaces from a group's own meshes lets those spaces be defined, documented
and shared on the project's terms, provided it copes with versatile mesh inputs. Second,
the scan data here (thin, fused and defective ant meshes) are an especially difficult
case, so human meshes collected under cleaner conditions, including public data sets,
form a subset of the complexity the pipeline is built to handle. Third, which joint
locations and body landmarks define a new model is decisive: graphics body models place
joints where skinning needs them, not at anatomical articulation centres, and their
simplified kinematic structure does not correspond to accurate joint locations~\cite{keller}.
The experiments below find the same sensitivity to definition (Sections~\ref{sec:skeleton}
and~\ref{sec:measurement}).

\subsection{Role and environment}

Under the supervision of Dr.\ Fabian Plum (author of scAnt~\cite{scant} and SMILify~\cite{smilify}), the work,
carried out on site from 1 July to 25 September 2026, extended an existing research code
base through Git and pull-request review. The differentiable fitter is implemented with
PyTorch3D~\cite{pytorch3d}, and large-scale experiments ran on GPU-enabled JURECA and RWTH
infrastructure, with experiment configurations, checkpoints and intermediate outputs
preserved for reproducibility.
